%%%%%%%%%%%%%%%%%%%%%%%%%%%%%%%%%%%%%%%%%%%%%%%%%%%%%%%%%%%%%%%%%%%%%%%%%%%%%%%%
%%
%% Abstract of the thesis in default document language
%%

% Drafted 2026-09-30 by an external LLM from prompts/PROMPT_final_chapters.md,
% audited and revised (see NOTES.md). Rewritten 2026-10-01 (author review v5
% #1-#2): more technical, results with numbers, no moving mesh, no comparison
% with another model. Numbers: Table 5.1 (mean +- sd over the five folds),
% Sections 5.3.2-5.3.4, 5.4.1, 5.4.3, 5.5.1 (growth-rate r).
% 2026-10-01: author decisions (a)-(c) applied (prompts/open.txt section 2):
% headline scoped (T-FEN above U-Net at both seeds, slightly above the
% stencil), global + local demoted to a one-seed observation, objective added.
% 2026-10-02 (author review v6 #2-#4): slot sentence starts from the slot,
% not the model list; results split into a framework paragraph and an
% operator paragraph (reach explained by the 21:1 spacing and the front
% advance, median ~12 columns ~ 0.07 mm per visit; k-NN one hop +-0.011 mm
% along a row, Fig. 4.2); spectral / local-path sentence dropped. "Plain
% mean squared error" = ANISOGNN_final_dilmean_plainmse (no mask weighting
% and no soft-Dice), Section 5.4.2.

\cleardoubleoddpage%  Make sure to start abstract on a new odd page
\begin{abstract*}
Geographic Atrophy (GA), the advanced atrophic form of age-related macular
degeneration, is a lesion of lost retinal tissue that grows slowly and
cannot be reversed. Forecasting how the lesion of an individual eye will
change is difficult because visits are irregularly spaced, each eye has
few of them, and cohorts are small.

This thesis asks what such a forecast needs: what kind of framework, and
what kind of operator inside it. Each visit is represented as an
eleven-channel state on a $49 \times 1024$ en-face grid from optical
coherence tomography (OCT): the GA mask and ten retinal layer depth maps.
A fixed framework advances this state by a residual update
$u + \Delta t\, f_\theta(u, \Delta t)$, conditioned on the time elapsed
between visits. A zero-initialised output starts every operator at exact
persistence, and all operators are trained with one pushforward curriculum
measured in elapsed days and one loss. Only the operator $f_\theta$
changes between experiments. It was filled in turn, at matched parameter
counts, with operators from four families: message-passing graph networks
on two neighbourhood graphs, a U-Net, Fourier Neural Operators with and
without a local path, and Finite Element Networks with and without a
learned transport term. They were compared by five-fold cross-validation
on 75 eyes of 51 patients from the Medical University of Vienna, using
the change-region Dice at one year.

Every operator with spatial context learned to forecast change; a
per-pixel model without it reproduced persistence exactly. The loss has
to emphasise the lesion mask: trained with a plain mean squared error,
the model learned no change at all, and a soft-Dice term added about
0.05. Neither fourth-order Runge--Kutta integration in place of a single
Euler step nor scaling the operators from about a quarter to four times
their size improved the forecast.

The three strongest operators come from three different families, a
Finite Element Network, a graph network and a U-Net, and reach a
change-region Dice of $0.54 \pm 0.04$, $0.53 \pm 0.05$ and
$0.50 \pm 0.06$ (mean $\pm$ standard deviation over five folds; zero for
persistence). What separates the operators is how far one update
reaches. The pixel spacing across B-scans is about 21 times larger than
along them, so a neighbourhood graph built on pixel indices reaches only
about 0.01\,mm along a B-scan per step, while the lesion front moves
about 0.07\,mm between visits. A dilated stencil that reaches about
0.12\,mm in both directions raises the Dice by 0.064 at an identical
parameter count, on all five folds. A learned transport term in the
Finite Element Network makes up for the same shortfall by a different
mechanism (+0.062 over the same network without it). The two operators
that carry these properties are the two highest.

The growth rate of an eye over its whole follow-up is predicted less well
than where the lesion changes (correlation 0.40 with the true rates). For
such forecasts, what matters is how far one update reaches compared with
how far the disease moves between visits, not the family of the operator.
\end{abstract*}

%%
%%%%%%%%%%%%%%%%%%%%%%%%%%%%%%%%%%%%%%%%%%%%%%%%%%%%%%%%%%%%%%%%%%%%%%%%%%%%%%%%


%%%%%%%%%%%%%%%%%%%%%%%%%%%%%%%%%%%%%%%%%%%%%%%%%%%%%%%%%%%%%%%%%%%%%%%%%%%%%%%%
%%
%% Abstract of the thesis in an additional language (German)
%%

\cleardoubleoddpage%  Make sure to start abstract on a new odd page
\begin{otherlanguage}{ngerman}
\begin{abstract*}
Die Geographische Atrophie (GA), die fortgeschrittene atrophe Form der
altersabhängigen Makuladegeneration, ist eine Läsion abgestorbenen
Netzhautgewebes, die langsam wächst und sich nicht zurückbildet.
Vorherzusagen, wie sich die Läsion eines einzelnen Auges verändern wird,
ist schwierig, weil die Kontrollen in unregelmäßigen Abständen
stattfinden, jedes Auge nur wenige davon hat und die Kohorten klein sind.

Diese Arbeit untersucht, was eine solche Vorhersage benötigt: welche Art
von Rahmen und welche Art von Operator darin. Jede Kontrolle wird als
Zustand mit elf Kanälen auf einem $49 \times 1024$ großen En-face-Gitter
der optischen Kohärenztomographie (OCT) dargestellt: der GA-Maske und
zehn Tiefenkarten retinaler Schichtgrenzen. Ein fester Rahmen schreibt
diesen Zustand mit einem residualen Update
$u + \Delta t\, f_\theta(u, \Delta t)$ fort, das von der zwischen zwei
Kontrollen verstrichenen Zeit abhängt. Eine mit null initialisierte
Ausgabe lässt jeden Operator bei exakter Persistenz beginnen, und alle
Operatoren werden mit einem in verstrichenen Tagen bemessenen
Pushforward-Curriculum und einer gemeinsamen Verlustfunktion trainiert.
Zwischen den Experimenten ändert sich nur der Operator $f_\theta$. Er
wurde nacheinander, bei abgeglichener Parameterzahl, mit Operatoren aus
vier Familien besetzt: Message-Passing-Graph-Netzwerken auf zwei
Nachbarschaftsgraphen, einem U-Net, Fourier Neural Operators mit und ohne
lokalen Pfad sowie Finite-Elemente-Netzwerken mit und ohne gelernten
Transportterm. Verglichen wurden sie mit fünffacher Kreuzvalidierung an
75 Augen von 51 Patienten der Medizinischen Universität Wien, anhand des
Dice-Koeffizienten der veränderten Region (Change-Region-Dice) nach einem
Jahr.

Jeder Operator mit räumlichem Kontext lernte, Veränderung vorherzusagen;
ein Modell, das jedes Pixel für sich betrachtet, reproduzierte exakt die
Persistenz. Die Verlustfunktion muss die Läsionsmaske betonen: Mit einem
einfachen mittleren quadratischen Fehler trainiert, lernte das Modell
keinerlei Veränderung, und ein Soft-Dice-Term brachte etwa 0,05. Weder eine
Runge--Kutta-Integration vierter Ordnung anstelle eines einzelnen
Euler-Schritts noch die Skalierung der Operatoren von etwa einem Viertel
bis zum Vierfachen ihrer Größe verbesserte die Vorhersage.

Die drei stärksten Operatoren stammen aus drei verschiedenen Familien, ein
Finite-Elemente-Netzwerk, ein Graph-Netzwerk und ein U-Net, und erreichen
einen Change-Region-Dice von $0{,}54 \pm 0{,}04$, $0{,}53 \pm 0{,}05$ und
$0{,}50 \pm 0{,}06$ (Mittelwert $\pm$ Standardabweichung über fünf Folds;
null für Persistenz). Was die Operatoren unterscheidet, ist, wie weit ein
einzelnes Update reicht. Der Pixelabstand quer zu den B-Scans ist etwa
21-mal größer als entlang der B-Scans, daher reicht ein auf Pixelindizes
gebauter Nachbarschaftsgraph pro Schritt entlang eines B-Scans nur etwa
0,01\,mm weit, während sich die Läsionsgrenze zwischen zwei Kontrollen um
etwa 0,07\,mm bewegt. Ein gedehnter Stencil, der in beide Richtungen etwa
0,12\,mm weit reicht, hebt den Dice bei identischer Parameterzahl um
0,064, auf allen fünf Folds. Ein gelernter Transportterm im
Finite-Elemente-Netzwerk gleicht dasselbe Defizit über einen anderen
Mechanismus aus (+0,062 gegenüber demselben Netzwerk ohne diesen Term).
Die beiden Operatoren mit diesen Eigenschaften sind die zwei stärksten.

Die Wachstumsrate eines Auges über die gesamte Beobachtungszeit wird
schlechter vorhergesagt als der Ort der Veränderung (Korrelation 0,40 mit
den tatsächlichen Raten). Für solche Vorhersagen zählt, wie weit ein
einzelnes Update reicht im Vergleich dazu, wie weit sich die Erkrankung
zwischen zwei Kontrollen bewegt, und nicht die Familie des Operators.
\end{abstract*}
\end{otherlanguage}

%%
%%%%%%%%%%%%%%%%%%%%%%%%%%%%%%%%%%%%%%%%%%%%%%%%%%%%%%%%%%%%%%%%%%%%%%%%%%%%%%%%
